\documentclass[10pt,a4paper]{amsart}
\usepackage{amsmath,amssymb,amsthm,graphicx}
\usepackage{verbatim}
\newenvironment{reference}{\comment}{\endcomment}
%\newenvironment{reference}{}{}
\newenvironment{slogan}{\comment}{\endcomment}
\newenvironment{history}{\comment}{\endcomment}
% For commutative diagrams we use Xy-pic
\usepackage[all]{xy}

% We use 2cell for 2-commutative diagrams.
\xyoption{2cell}
% Theorem environments.
%
\theoremstyle{plain}
\newtheorem{theorem}[subsection]{Theorem}
\newtheorem{proposition}[subsection]{Proposition}
\newtheorem{lemma}[subsection]{Lemma}

\theoremstyle{definition}
\newtheorem{definition}[subsection]{Definition}
\newtheorem{example}[subsection]{Example}
\newtheorem{exercise}[subsection]{Exercise}
\newtheorem{situation}[subsection]{Situation}

\theoremstyle{remark}
\newtheorem{remark}[subsection]{Remark}
\newtheorem{remarks}[subsection]{Remarks}

\numberwithin{equation}{subsection}

% Macros
%
\def\lim{\mathop{\mathrm{lim}}\nolimits}
\def\colim{\mathop{\mathrm{colim}}\nolimits}
\def\Spec{\mathop{\mathrm{Spec}}}
\def\Hom{\mathop{\mathrm{Hom}}\nolimits}
\def\Ext{\mathop{\mathrm{Ext}}\nolimits}
\def\SheafHom{\mathop{\mathcal{H}\!\mathit{om}}\nolimits}
\def\SheafExt{\mathop{\mathcal{E}\!\mathit{xt}}\nolimits}
\def\Sch{\mathit{Sch}}
\def\Mor{\mathop{\mathrm{Mor}}\nolimits}
\def\Ob{\mathop{\mathrm{Ob}}\nolimits}
\def\Sh{\mathop{\mathit{Sh}}\nolimits}
\def\NL{\mathop{N\!L}\nolimits}
\def\CH{\mathop{\mathrm{CH}}\nolimits}
\def\proetale{{pro\text{-}\acute{e}tale}}
\def\etale{{\acute{e}tale}}
\def\QCoh{\mathit{QCoh}}
\def\Ker{\mathop{\mathrm{Ker}}}
\def\Im{\mathop{\mathrm{Im}}}
\def\Coker{\mathop{\mathrm{Coker}}}
\def\Coim{\mathop{\mathrm{Coim}}}

% Boxtimes
%
\DeclareMathSymbol{\boxtimes}{\mathbin}{AMSa}{"02}

%
% Macros for moduli stacks/spaces
%
\def\QCohstack{\mathcal{QC}\!\mathit{oh}}
\def\Cohstack{\mathcal{C}\!\mathit{oh}}
\def\Spacesstack{\mathcal{S}\!\mathit{paces}}
\def\Quotfunctor{\mathrm{Quot}}
\def\Hilbfunctor{\mathrm{Hilb}}
\def\Curvesstack{\mathcal{C}\!\mathit{urves}}
\def\Polarizedstack{\mathcal{P}\!\mathit{olarized}}
\def\Complexesstack{\mathcal{C}\!\mathit{omplexes}}
% \Pic is the operator that assigns to X its picard group, usage \Pic(X)
% \Picardstack_{X/B} denotes the Picard stack of X over B
% \Picardfunctor_{X/B} denotes the Picard functor of X over B
\def\Pic{\mathop{\mathrm{Pic}}\nolimits}
\def\Picardstack{\mathcal{P}\!\mathit{ic}}
\def\Picardfunctor{\mathrm{Pic}}
\def\Deformationcategory{\mathcal{D}\!\mathit{ef}}
\usepackage{cleveref}
\crefname{originallemma}{original lemma}{original lemmas}
\Crefname{originallemma}{Original lemma}{Original lemmas}
\makeatletter
\newlabel{lemma-inverse-limit-spectral-spaces-constructible}{{lemma-inverse-limit-spectral-spaces-constructible}{source}{}{sourcepointer-1}{}}
\newlabel{lemma-inverse-limit-spectral-spaces-constructible@cref}{{[originallemma][0][]lemma-inverse-limit-spectral-spaces-constructible}{[1][0][]source}{}{}{}}
\makeatother
\crefname{originallemma}{original lemma}{original lemmas}
\Crefname{originallemma}{Original lemma}{Original lemmas}
\makeatletter
\newlabel{lemma-inverse-limit-spectral-spaces-equal}{{lemma-inverse-limit-spectral-spaces-equal}{source}{}{sourcepointer-2}{}}
\newlabel{lemma-inverse-limit-spectral-spaces-equal@cref}{{[originallemma][0][]lemma-inverse-limit-spectral-spaces-equal}{[1][0][]source}{}{}{}}
\makeatother
\title{Quasi-compact opens of an inverse limit of spectral spaces descend to a finite stage}
\author{The Stacks Project Authors}
\date{stacks-project:topology.tex:local-snapshot-20261012; GNU FDL 1.2 or later}
\begin{document}
\maketitle
\noindent\textit{Editorial note.} All delivered mathematical text is the author's own, copied from the pinned Stacks Project chapter source sources/stacks-topology/pinned-original-source.tex; no statement or proof step was rewritten, reordered or removed.\par\medskip
\noindent\textit{Editorial note.} Preamble provenance: the Stacks Project keeps one single shared preamble for every chapter; this card takes it directly from the staged sources/stacks-topology/preamble.tex as three verbatim spans (lines 8-12, 14-18 and 154-217), so no derived file was written and nothing was added or redefined. The wrapper class and the base packages are supplied by the assembler.\par\medskip
\noindent\textit{Editorial note.} Bibliography provenance: neither the counted statement nor its proof contains a citation (this lane rescans the whole span before writing the card), so the card carries no bibliography and no entry was reconstructed.\par\medskip
\noindent\textit{Editorial note.} Selection provenance: this card was opened under the file-based control contract (CONTROL.md), with the mechanical pre-card gate run on its statement and proof interval before the card was written.\par\medskip
\noindent\textit{Allowed standard prerequisites (editorial).}\begin{itemize}
\item The statement that constructible subsets of the limit come from a finite stage. Invoked inside the counted proof (source lines 5211-5213). Descends the subset to a stage.
\item The statement on differences of pullbacks in such a limit. Invoked inside the counted proof (source lines 5214-5216). Provides the nonemptiness criterion for the differences appearing in the statement.
\end{itemize}\medskip
\ifdefined\hypertarget\hypertarget{sourcepointer-1}{}\fi
\noindent\textit{Source reference.} Original lemma lemma-inverse-limit-spectral-spaces-constructible: sources/stacks-topology/pinned-original-source.tex, source lines 5079--5079. Navigation pointer only: the counted text refers to the Stacks item lemma-inverse-limit-spectral-spaces-constructible, which is an existing labeled result of this same chapter and is not part of this card. The pointer carries the target's own label as its number and its definition line as its span; its content is not reproduced here.\par\medskip
\ifdefined\hypertarget\hypertarget{sourcepointer-2}{}\fi
\noindent\textit{Source reference.} Original lemma lemma-inverse-limit-spectral-spaces-equal: sources/stacks-topology/pinned-original-source.tex, source lines 5053--5053. Navigation pointer only: the counted text refers to the Stacks item lemma-inverse-limit-spectral-spaces-equal, which is an existing labeled result of this same chapter and is not part of this card. The pointer carries the target's own label as its number and its definition line as its span; its content is not reproduced here.\par\medskip
\begin{lemma}
\label{lemma-descend-opens}
Let $\mathcal{I}$ be a cofiltered index category.
Let $i \mapsto X_i$ be a diagram of spectral spaces such
that for $a : j \to i$ in $\mathcal{I}$ the corresponding map
$f_a : X_j \to X_i$ is spectral. Set $X = \lim X_i$ and denote
$p_i : X \to X_i$ the projection.
\begin{enumerate}
\item Given any quasi-compact open $U \subset X$
there exists an $i \in \Ob(\mathcal{I})$ and a quasi-compact open
$U_i \subset X_i$ such that $p_i^{-1}(U_i) = U$.
\item Given $U_i \subset X_i$ and $U_j \subset X_j$
quasi-compact opens such that $p_i^{-1}(U_i) \subset p_j^{-1}(U_j)$
there exist $k \in \Ob(\mathcal{I})$ and morphisms
$a : k \to i$ and $b : k \to j$ such that $f_a^{-1}(U_i) \subset f_b^{-1}(U_j)$.
\item If $U_i, U_{1, i}, \ldots, U_{n, i} \subset X_i$ are quasi-compact
opens and
$p_i^{-1}(U_i) = p_i^{-1}(U_{1, i}) \cup \ldots \cup p_i^{-1}(U_{n, i})$
then
$f_a^{-1}(U_i) = f_a^{-1}(U_{1, i}) \cup \ldots \cup f_a^{-1}(U_{n, i})$
for some morphism $a : j \to i$ in $\mathcal{I}$.
\item Same statement as in (3) but for intersections.
\end{enumerate}
\end{lemma}

\begin{proof}
Part (1) is a special case of
Lemma \ref{lemma-inverse-limit-spectral-spaces-constructible}.
Part (2) is a special case of
Lemma \ref{lemma-inverse-limit-spectral-spaces-equal}
as quasi-compact opens are both open and closed in the constructible
topology. Parts (3) and (4) follow formally from (1) and (2)
and the fact that taking inverse images of subsets commutes with
taking unions and intersections.
\end{proof}

\end{document}
